\section{Experimental Setup}
\label{sec:setup}

\subsection{Datasets}

\begin{table}[h]
\caption{Dataset summary.}
\label{tab:datasets}
\begin{tabular}{lcccc}
\toprule
Dataset & Domain & Users & Items & Demographics \\
\midrule
MovieLens-1M~\cite{harper2016movielens} & movie & 6{,}031 & 3{,}883 & gender $\times$ age \\
LastFM-1K~\cite{celma2010longtail} & music track & 942 & 272{,}769 & gender $\times$ age \\
\bottomrule
\end{tabular}
\end{table}

ML-1M provides explicit 1--5 ratings with complete demographic fields. LFM-1K provides per-track listening counts; demographic fields are partial (gender $89\%$, age $29\%$); we treat missing as \textsf{Unknown} slice.

For LFM-1K, item granularity is \textbf{track-level} (not artist-level): each item key is \texttt{Artist|||Track}. This is a more demanding setup than artist-level due to extreme sparsity (max popularity count of 224 over 992 users).

\subsection{LLM Recommender}

We use \textbf{GPT-4.1-mini} as the underlying LLM recommender. The same prompt format is used for both datasets, with domain-specific templates (movie / music\_track) injected at render time. Demographic descriptors are added per scenario:
\begin{itemize}
\item $g = \emptyset$: no descriptor (neutral)
\item $g = \textsf{gender}$: e.g., \emph{``the user is female''}
\item $g = \textsf{age}$: e.g., \emph{``the user is in the young adult bracket''}
\item $g = \textsf{cross}$: combined
\end{itemize}

For each (user, scenario) we issue one LLM call, parse the numbered list, and obtain $N=30$ candidate items. Cache hits are $100\%$ on subsequent reranking-config experiments, removing API cost from ablations.

\subsection{Reranker Configuration}

Default v2 configuration (all formulas in Section~\ref{sec:method}):

\begin{table}[h]
\caption{Default reranker hyperparameters.}
\label{tab:hyperparams}
\begin{tabular}{ll}
\toprule
Parameter & Value \\
\midrule
$K$ (output length) & 15 (ML-1M), 20 (LFM-1K) \\
$N$ (candidate pool) & 30 \\
$L$ (history limit fed to LLM) & 25 \\
$\text{base}$ (preference floor) & 0.2 \\
$r$ (decay rate) & 0.9 \\
$\alpha_{\text{base}}, \alpha_{\min}, \alpha_{\max}$ & 0.7, 0.5, 0.9 \\
$\alpha_{\text{gain}}$ & 0.2 \\
$B$-transform & \texttt{rank} \\
$K'$ (head protection) & 5 \\
$\rho$ (rank-bonus weight) & 0.2 \\
$\theta$ (APLT threshold) & 0.2 \\
\bottomrule
\end{tabular}
\end{table}

\subsection{Statistical Inference}

Per-user paired differences are computed and analyzed via:
\begin{itemize}
\item Paired $t$-statistic $t = \bar{d} / (s_d / \sqrt{n})$
\item Two-sided $p$-value via normal approximation (justified by $n \geq 600$)
\item $95\%$ confidence interval $\bar{d} \pm 1.96 \cdot \text{SE}(\bar{d})$
\item Cohen's $d = \bar{d} / s_d$ for effect size
\end{itemize}

Bonferroni correction is applied for multiple comparisons across (scenarios $\times$ metrics).

\subsection{Reproducibility}

All code, configurations, and intermediate artifacts (popularity tables, adapter pickle caches, LLM response caches) are released. Random seed is fixed at 42 for sampling, splitting, and prompt rendering. End-to-end pipeline reproducibility is ensured through three commands documented in the supplementary repository.
